\scanpage{173}
\mspage{89}
% archive sheet number: 172
% lines 1-8
Ce lemme est une tautologie: si $F$ désigne la limite inductive dans $\widehat C$, il est clair \inserted{par définition} (en prenant un foncteur $I_0\mathrel{\rule{1.5em}{0.4pt}}I$ cofinal, avec $I_0$ \uline{ordonné} grand devant $\pi$, de sorte que la $\varinjlim$ peut être prise suivant $I_0$) que $F$ est dans $\operatorname{Ind}'_{\pi}(C)$, et donc c'est à fortiori une $\varinjlim$ dans $\operatorname{Ind}'_{\pi}C$.

% lines 9-15
Pour vérifier l'existence des $\varinjlim$ de cardinal $\leq\pi$ dans $\operatorname{Ind}_{\unsure{\pi|\beta}}C$, dans le cas où $C$ lui-$=$ est stable par \unsure{petites|telles} limites, il suffit donc, grâce au corollaire 4.4.8, de prouver la cofinalité des foncteurs (4.4.6.3). Cela va résulter du \deleted{\unsure{stable par $\varinjlim$ de cardinal $\leq\pi$}}

% lines 16-26
\uline{Lemme} 4.4.\unsure{9|8}. Soient \deleted{\unsure{$\pi$}}, \inserted{$\pi$ un cardinal infini,} $C$ une petite catégorie, $J$ une petite catégorie. \deleted{Supposons $C$ stable par sommes de cardinal $\leq\pi$, et} \deleted{telle} \inserted{que} pour tous $j,j'\in\operatorname{Ob}J$, on ait $\operatorname{card}\operatorname{Hom}(j,j')\leq\pi$. \deleted{Soit alors} Dans ces conditions, \deleted{les} \inserted{pour tout $j\in\operatorname{Ob}J$,} \unsure{le} foncteur (4.4.6.3) \deleted{\unsure{est un}} \uline{admet un adjoint à gauche}, donc à fortiori il est cofinal (car les catégories colocalisées en dessous des $C_{/\varphi(j)}$ ont un objet initial, à fortiori elles sont $0$-connexes).
% margin note written slanted, crossed out
\margin{\deleted{\unsure{Adjonctions}}}
